\documentclass[10pt]{article}

\usepackage[utf8]{inputenc}

\usepackage[T1]{fontenc}

\usepackage{amsmath}

\usepackage{amsfonts}

\usepackage{amssymb}

\usepackage[version=4]{mhchem}

\usepackage{stmaryrd}

\usepackage{bbold}

\usepackage{graphicx}

\usepackage[export]{adjustbox}

\graphicspath{ {./images/} }



\begin{document}

резонансной частоте оптически разрешенного перехода ее атомов, молекул или примесных ионов. Как впервые обнаружено в [5], сигнал стимулированного Ф. э. при нек-рых условиях может повторять форму одного из возбуждающих импульсов. С другой стороны, временной интервал между первыми двумя (информационными) и третьим (считывающим) импульсами может быть весьма велик. Так, в эксперименте (см. [6]) его удалось довести до 13 часов.



Лит.: [1] Ко п в ил л ем У.Х., Нагибаров В. Р., «Физика металлов и металловедение», 1963, т. 15, № 2, с. 313-15; [2] Kurnit N.A., AbellaJ.D., HartmannS.R., «Phys. Rev. Lett.», 1964, v. 13, № 19, p. 567-68; [3] М аныкин Э. А., С а м а р це в В. В., Оптическая эхо-спектроскопия, М., 1984; [4] Е в с е ев И. В., Ермаченко В. М., в кн.: IV Междунар. симпозиум по избранным проблемам статистической механики, Дубна, 1988, с. 107-14; [5] Зуйков В. А., Самарцев В. В., Усманов Р.Г., «Письма в ЖЭТФ», 1980, т. 32, в. 4, с. 293-97; [6] B abb itt W. R., Mossberg T. W., «Opt. Comm.», 1988, v. 65, № 3, p. 185-88.



И. В. Евсеев.



ФОтООтС4ETOB СТATИСТИКА - результат обработки данных квантовооптического эксперимента и сопоставления с модельными представлениями, вытекающими из $\mathrm{cma}-$ тистической теории отсчетов излучения. Напр., пусть оптич. излучение с комплексной амплитудой $\xi(t)$ является суперпозицией случайной амплитуды $\alpha(t)$ и детерминированной амплитуды $\beta(t)$. Если $\alpha(t)$ является решением стохастич. дифференциального уравнения



\[

d \alpha(t)=-v \alpha(t) d t+\sqrt{\sigma} d w(t)

\]



$v, \sigma=$ const, $w(t)$ - стандартный комплекснозначный винеровский процесс, то производящая функция $Q(\lambda)$ случайного числа $N$ фотоотсчетов



\[

Q(\lambda)=\langle\exp (-\lambda N)\rangle

\]



регистрируемых на интервале $[0, T]$, равна



\[

\begin{aligned}

& Q(\lambda)=\left\langle\exp \left\{-\lambda \int_{0}^{T}|\alpha(t)+\beta(t)|^{2} d t\right\}\right\rangle= \\

& =4 r v e^{v T}\left(r_{+}^{2} e^{r T}-r_{-}^{2} e^{-r T}\right)^{-1} \times \\

& \times \exp \left\{-\lambda \int_{0}^{T}|\beta(t)|^{2} d t-\frac{2 \lambda^{2} \sigma r^{-1}}{r_{+}^{2} e^{r T_{-} r_{-}^{2} e^{-r T}}} \times\right. \\

& \left.\times \int_{0}^{T} d t \int_{t}^{T} d \tau\left(r_{+} e^{r t}+r_{-} e^{-r t}\right)\left(r_{+} e^{r(T-\tau)}+r_{-} e^{-r(T-\tau)}\right)\right\},

\end{aligned}

\]



где $r=\left(v^{2}+2 \lambda \sigma\right)^{1 / 2}, r_{ \pm}=r \pm v$, а распределение вероятностей $P(n)=\mathrm{P}\{N=n\}$ случайного числа $N$ фотоотсчетов на интервале $[0, T]$ равно



\[

P(n)=\left.\frac{(-1)^{n}}{n !} \frac{d^{n}}{d \lambda^{n}} Q(\lambda)\right|_{\lambda=1} ;

\]



оно используется при анализе данных эксперимента.



Лит.: [1] Клауде р Дж., Сударшан Э., Основы квантовой оптики, пер. с англ., М., 1970; [2] Мазман иш в или А. С., Континуальное интегрирование как метод решения физических задач, Киев, 1987. А. С. Мазманишвили. ФРАГМЕНА - ЛИНДЕЛЕФА ПРИНЧИП - семейство высказываний, распространяющих максимума модуля принцип для голоморфной функции на случай неограниченной области. Типичное утверждение такого рода: пусть функция $f(z)$ голоморфна в угле раствора $\pi / \alpha$ и непрерывна вплоть до его сторон. Если $|f(z)| \leqslant M$ на сторонах угла, то либо $|f(z)| \leqslant M$ всюду внутри угла, либо



\[

\lim _{\rho \rightarrow \infty} \rho^{-\alpha} \ln \max _{|z|=\rho}|f(z)|>0 .

\]



Имеются обобщения Ф. - Л. п. для решений и субрешений дифференциальных уравнений эллиптич. типа, для пространственных отображений с ограниченным искажением.



В. М. Миклюков. ФРАГМЕНА - ЛИНДЕЛЕФА ТЕОРЕМА - обобщение максимума модуля принципа аналитических функций на случай функций, априори заданных как неограниченные; впервые в простейшей форме дано Э. Фрагменом и Э. Линделефом (E. Phragmen, E. Lindelöf, 1908). Пусть $f(z)-$ peгулярная аналитич. функция комплексного переменного z в области $D$ на плоскости $\mathbb{C}$ с границей Г. Говорят, что $f(z)$ не превосходит по модулю числа $M$ в граничной точке $\zeta \in \Gamma$, если



\[

\lim _{z \rightarrow \zeta} \sup _{z \in D}|f(z)| \leqslant M

\]



то есть если для любого $\varepsilon>0$ найдется круг $\Delta$ (зависящий от $\zeta$ и $\varepsilon)$ с центром $\zeta$ такой, что $|f(z)|<M+\varepsilon$ при $z \in D \cap \Delta$. Основное содержание результатов Э. Фрагмена и Э. Линделефа, в несколько модернизированной форме, заключается в следующих двух положениях, последовательно расширяющих принцип максимума модуля.



I. Если регулярная аналитич. функция $f(z)$ всюду на границе $\Gamma$ не превосходит по модулю числа $M$, то $|f(z)| \leqslant M$ всюду в области $D$. Это положение иногда называют Фрагмена - Линделефа принципом. Оно расширяет принцип максимума модуля на функции, о граничном поведении к-рых имеется лишь частичная информация.



II. Пусть регулярная аналитич. функция $f(z)$ не превосходит по модулю числа $M$ во всех точках границы $\Gamma$, не принадлежаших нек-рому множеству $E \subset \Gamma$. Пусть, кроме того, существует функция $\omega(z)$ со следующими свойствами: 1) $\omega(z)$ регулярна в $D$, 2) $|\omega(z)|<1$ в $D$, 3) $\omega(z) \neq 0$ в $D$, 4) для любого $\sigma>0$ функция $|\omega(z)|^{\sigma}|f(z)|$ не превосходит числа $M$ в каждой точке $\zeta \in E$. При этих условиях $|f(z)| \leqslant M$ всюду в $D$.



Формулировки I и II остаются в силе для голоморфной функции $f(z), z=\left(z_{1}, \ldots, z_{n}\right)$, заданной в области $D$ комплексного пространства $\mathbb{C}^{n}, n \geqslant 1$. Многие работы посвяшены получению результатов типа Ф. - Л. т. для решений дифференциальных уравнений с частными производными и систем уравнений эллиптич. типа.



Лuт.: [1] Е в г р а фо в М.А., Аналитические функции, 2 изд., М., 1968. [1] Е в г ра фо в М. А., Аналитические функи. Соломенцев.

Е. Д. ФРЕДГОЛЬМА АЛЬТЕРНАТИВА - см. Интегральное уравнение.



ФРЕДГОЛЬМА ИНТЕГРАЛЬНОЕ УРАВНЕНИЕ см. Интегральное уравнение.



ФРЕДГОЛЬМОВ OПEPATOP - см. Интегральное уравнение, Интегральный оператор.



ФРЕДГОЛЬМОВО ЯДРО - см. Интегральный оператор. ФРЕЛИХА МОДЕЛЬ - модель, предложенная Х. Фрелихом [1] для описания блоховских электронов, взаимодействующих с фононами в металлах, гамильтониан которой



\[

\begin{gathered}

H=\sum_{\vec{k}, \sigma}\left(\varepsilon_{\vec{k}}-\mu\right) a_{\vec{k}, \sigma}^{+} a_{\vec{k}, \sigma}+\sum_{\vec{q}} \hbar \omega_{q} b_{\vec{q}}^{+} b_{\vec{q}}+ \\

+V^{-1 / 2} \sum_{\vec{k}, \sigma, \vec{q}} g(q)\left[a_{\vec{k}+\vec{q}, \sigma}^{+} a_{\vec{k}, \sigma} b_{\vec{q}}+a_{\vec{k}, \sigma}^{+} a_{\vec{k}+\vec{q}, \sigma} b_{\vec{q}}^{+}\right] .

\end{gathered}

\]



Канонич. преобразованием во втором порядке теории возмущений в гармонич. приближении для фононов было получено (см. [1]) эффективное взаимодействие между электронами, не зависящее от фононных амплитуд, с перенормировкой фононного спектра и константы связи:



\[

\begin{aligned}

& H_{I}=\frac{1}{2} \sum_{\vec{k}, \overrightarrow{k^{\prime}}, \sigma, \sigma^{\prime}, \vec{q}} \frac{2 \hbar \omega_{q}\left|M_{q}\right|^{2}}{\left(\varepsilon_{\vec{k}}-\varepsilon_{\vec{k}+\vec{q}}\right)^{2}-\left(h \omega_{q}\right)^{2}} \times \\

& \times a_{\overrightarrow{k^{\prime}}}^{+}-\vec{q}, \sigma^{\prime} a_{\overrightarrow{k^{\prime}}, \sigma^{\prime}} a_{\vec{k}+\vec{q}, \sigma}^{+} a_{\vec{k}, \sigma} .

\end{aligned}

\]



\begin{center}

\includegraphics[max width=\textwidth]{2023_07_08_d36ded0585bb6cd4446bg-1}

\end{center}



Jum.: [1] Frohlich H., «Proc. Roy. Soc. (London)», 1952, A 215, p. 291; [2] Bardeen J., Pines D., «Phys. Rev.», 1955, v. 99, p. 1140.





\end{document}